% Source: 03 - Wed Sep 30 2026.pdf, page 7. Content remains in source order.
\sourcepage{03}{7}
$\Longleftrightarrow$ means that we need to prove two implications:
\begin{itemize}
\item strings in $L_1$ are mapped to strings in $L_2$ ($\Rightarrow$);
\item strings not in $L_1$ are mapped to strings not in $L_2$ ($\Leftarrow\ \equiv\ \not\Rightarrow$: COUNTERPOSITIVE).
\end{itemize}

\textbf{Remark:} $f$ need not be neither injective or surjective.
\sourcefigure[0.9]{assets/03/page-007-reduction-map.png}

Poly-time reducibility captures the concept of ``lesser difficulty'' of $L_1$ with respect to $L_2$ in the following sense: $L_1\polyred L_2$. 

 $L_1$ is at most as difficult as $L_2$ but it cannot be more difficult than $L_2$ (otherwise, we would not be able to reduce $L_1$ to $L_2$). This is a way to put an order between problems. 

 We can now study $L_2$ to get information about $L_1$. If  $L_2$ can be solved in polynomial time, then $L_1$ must be solvable in polynomial time as well. If $L_2$ is not solvable in polynomial time, then we cannot conclude anything about $L_1$.

